%======================================================================
% INTRODUCCIÓN
%======================================================================

\chapter{Introducción}

Sandra UI Consola es la capa de presentación del ecosistema Sandra, proporcionando una interfaz administrativa robusta, segura y de alto rendimiento para la gestión integral del sistema. Esta consola representa el nodo de control de un ecosistema distribuido diseñado para la resiliencia y la observabilidad completa.

\section{Acerca de Sandra UI}

En el contexto de la orientación a objetos, la lógica del negocio es aquella funcionalidad ofrecida por el software. El software se comunica de manera amigable con el usuario a partir de la interfaz, pero el procesamiento de los datos capturados como entrada y la posterior entrega de resultados al usuario por medio de la interfaz es conocido como la Lógica de Negocio.

Sandra UI Consola se posiciona como una solución Enterprise de alta gama, integrando capacidades de SRE (Site Reliability Engineering) y observabilidad avanzada para garantizar la continuidad operativa. La consola proporciona una experiencia de usuario premium basada en principios de diseño moderno, mientras mantiene los más altos estándares de seguridad y rendimiento.

Los principios de SRE permiten combinar la velocidad de desarrollo con la fiabilidad operacional, reduciendo el tiempo medio de recuperación (MTTR) y aumentando la disponibilidad del sistema.

\section{Ecosistema Tecnológico}

Sandra UI Consola está construida sobre un ecosistema tecnológico de alta gama que garantiza rendimiento, seguridad y escalabilidad.

\subsection{Backend Core}

El núcleo del backend utiliza tecnologías de alto rendimiento:

\begin{itemize}
\item \textbf{Go}: Lenguaje de programación diseñado para maximizar la eficiencia en el manejo de memoria y concurrencia, proporcionando tiempos de respuesta ultrarrápidos.
\item \textbf{Rust}: Utilizado para componentes críticos que requieren seguridad de memoria y procesamiento de alto rendimiento.
\end{itemize}

La combinación de Go y Rust permite aprovechar las fortalezas de ambos lenguajes: la facilidad de desarrollo y concurrencia de Go, junto con la seguridad de memoria y rendimiento de Rust.

\subsection{Comunicación}

La consola implementa múltiples protocolos de comunicación:

\begin{itemize}
\item \textbf{gRPC con Protocol Buffers}: Serialización de datos de alta velocidad y baja latencia para streaming de información en tiempo real.
\item \textbf{REST API}: Endpoints estándar para operaciones CRUD y consultas.
\item \textbf{WebSockets}: Comunicación bidireccional para actualizaciones en tiempo real sin polling.
\item \textbf{AES-256}: Cifrado de punto a punto para datos en reposo.
\item \textbf{RSA/TLS 1.3}: Protocolo de seguridad para datos en tránsito.
\end{itemize}

La combinación de gRPC y WebSockets permite implementar patrones de comunicación tanto síncronos como asíncronos, optimizando el uso de recursos de red.

\subsection{Base de Datos}

El sistema soporta un modelo de almacenamiento híbrido:

\begin{itemize}
\item \textbf{MongoDB}: Base de datos NoSQL orientada a documentos, optimizada para grandes volúmenes de datos y flexibilidad de esquema.
\item \textbf{Motores Relacionales}: Soporte para bases de datos SQL cuando se requiere integridad transaccional.
\end{itemize}

\section{Requisitos del Sistema}

Para utilizar Sandra UI Consola se requieren los siguientes componentes:

\begin{itemize}
    \item Navegador web moderno (Chrome, Firefox, Edge o Safari)
    \item Conexión a una instancia de Sandra Server
    \item Credenciales de acceso válidas
    \item Resolución de pantalla mínima de 1280x720
\end{itemize}

\clearpage
